% Current v15 frontier. Stable v13/v14 shape IDs; no existence assertion.
\subsection{残る十三型と，次の具体的な入口}
\label{ns15:support-agenda}

高分母 support の必要条件を，$K_5$ の全 $2^{10}=1024$ 個の標識付き部分 graph に適用する．
木・単純閉路成分の除去後は二十型・六百八十六標識，
二価二頂点を持つ三角形の除去後は十六型・五百二十一標識，
葉辺付き四角形の除去後は十五型・四百六十一標識であった．
v15 では diamond $G_{4,2}$ の三十標識と $K_{2,3}=G_{5,6}$ の十標識を排除したので，
現在の候補は\emph{十三型・四百二十一標識}である．
以下の $G_{n,j}$ の番号は以前と同じ識別子を保ち，欠番を詰めない．
孤立頂点を省き，$ij$ は無向辺 $\{i,j\}$ を表す．
$\mathbf d$ は降順の次数列，$\beta_1=|E|-n+1$ は cycle rank，$L$ は葉の数，
$N$ は $K_5$ 上の標識付き個数である．

\begin{center}
\footnotesize
\setlength{\tabcolsep}{3pt}
\begin{tabular}{crclrrr}
型 & $|E|$ & $\mathbf d$ & $E$ & $\beta_1$ & $L$ & $N$\\\hline
$G_{4,3}$ & 6 & $(3333)$ & $\{01,02,03,12,13,23\}$ & 3 & 0 & 5 \\
$G_{5,3}$ & 5 & $(33211)$ & $\{01,02,03,12,14\}$ & 1 & 2 & 60 \\
$G_{5,5}$ & 6 & $(33222)$ & $\{01,02,03,12,14,34\}$ & 2 & 0 & 60 \\
$G_{5,7}$ & 6 & $(33321)$ & $\{01,02,03,12,13,24\}$ & 2 & 1 & 60 \\
$G_{5,9}$ & 6 & $(43221)$ & $\{01,02,03,04,12,13\}$ & 2 & 1 & 60 \\
$G_{5,10}$ & 7 & $(33332)$ & $\{01,02,03,12,13,24,34\}$ & 3 & 0 & 30 \\
$G_{5,11}$ & 7 & $(43322)$ & $\{01,02,03,04,12,13,24\}$ & 3 & 0 & 60 \\
$G_{5,12}$ & 7 & $(43331)$ & $\{01,02,03,04,12,13,23\}$ & 3 & 1 & 20 \\
$G_{5,13}$ & 7 & $(44222)$ & $\{01,02,03,04,12,13,14\}$ & 3 & 0 & 10 \\
$G_{5,14}$ & 8 & $(43333)$ & $K_5-\{14,23\}$ & 4 & 0 & 15 \\
$G_{5,15}$ & 8 & $(44332)$ & $K_5-\{24,34\}$ & 4 & 0 & 30 \\
$G_{5,16}$ & 9 & $(44433)$ & $K_5-\{34\}$ & 5 & 0 & 10 \\
$G_{5,17}$ & 10 & $(44444)$ & $K_5$ & 6 & 0 & 1 \\
\end{tabular}
\end{center}

同じ次数列でも異なる graph があるため辺集合を併記した．
代表は使用頂点の全置換の下で辞書式最小の辺リストを選ぶ．
全ての標識付き graph と自己同型群を
\nolinkurl{verify_support_agenda_v15.py} で再検算している．
表は \nolinkurl{support_agenda_verification_v15.json} の \texttt{shapes} から再生成できる．
この十三型は必要条件を通過した候補であり，整数接着の存在や atlas の同型類の数ではない．

\paragraph{現在の構造的な結論．}
非空の support は必ず三角形を含む．triangle-free で非自明な閉路を持つ候補は，
$C_4$，$C_5$，葉辺付き $C_4$，$K_{2,3}$ に尽くされ，いずれも排除された．
従って triangle-free support では全ての分母が一である．
四頂点を使う非空 support は $G_{4,3}=K_4$ だけに限られる．
さらに，二価 core が負の枝なら，その高分母近傍間も高辺でなければならない．
この最後の条件には対向 $D$ の相殺は不要である．

\paragraph{第一候補：異なる頂点に二本の葉を持つ三角形 $G_{5,3}$．}
これは三角形 $012$ と葉辺 $03,14$ からなる唯一の残存 unicyclic 型である．
葉の係数は中心 $0,1$ にそれぞれ $D\geq q$ を要求する．
二価 core $2$ の二高分母標準形と，二つの葉 $3,4$ の低 rows を同時に課す．
canonical 高 rays を $A=0,B=1$ とすると，$C=3,D=4$ の対応は排除され，
$C=4,D=3$ という交差対応だけが残る．さらに中央 core の負の枝も排除された．
$S=T=w\geq5$ は一般の不等式で，$w\leq4$ は証明で有限化した七つの局所模型と
四百八十五の葉の組の完全復元で処理する．
従って残る中央 core は
\[
 S>r,\qquad rT>(r+1)S,\qquad r\leq4\ \text{または}\ z=1
\]
を満たす正の枝である．これは $G_{5,3}$ 全体の排除ではなく，この型は未解決として残す．
次はこの二つの正の枝の範囲で，両葉の marked entries と node 数を中心へ伝え，
整数 clutching の剰余と対向 $D$ の要求を同時に整理する．

\paragraph{第二候補：四頂点の完全 graph $G_{4,3}=K_4$．}
第五頂点は全低分母で，残る四つの core はいずれも三高分母である．
diamond の証明で得た，全低 core に接する四本の $D$ の総和恒等式と，
共通低 puncture における符号付き約数の三択を再利用する．
ただし $K_4$ では四中心間の全辺が高分母なので，diamond の唯一の低辺が担った
低低 entry の同期はそのまま使えない．中心の低 row の node 条件と
各高辺の例外の選択を組み合わせることが具体的な入口となる．
三高分母の matching 性を無条件の既成事実として使用しない．

\paragraph{第三候補：二価 core の多い $G_{5,5}$ と $G_{5,10}$．}
$G_{5,5}$ の二価頂点 $3,4$ は，高分母近傍間に低辺を持つので，
どちらも二高分母標準形の正の枝に限られる．
頂点 $2$ の高分母近傍 $0,1$ は高辺で結ばれるため，この頂点では負の枝も残す．
二つの正の枝の低辺 node 条件と，三角形側の一つの二高分母 core を接続する．
$G_{5,10}$ では二価頂点 $4$ だけが同じ理由で正の枝に限られる．
これらは旧 $K_{2,3}$ 必要条件系の単純な転用ではなく，新しい近傍間低辺の補題が与える制約である．

\paragraph{他の型と，局所理論の保留事項．}
残る十三型には，葉・二価頂点・次の頂点という長さ二の枝がない．
完全葉伝播は再利用可能な補題として保持するが，この枝形を再探索する必要はない．
$G_{5,14}=K_{2,2,1}$ では既存の dense mod 五の二模型の整数 lifting を検討できる．
高分母 support $H=\{ij:q_{ij}>1\}$ と素数 support $G_\ell=\{ij:\ell\mid q_{ij}\}$ を区別する．
三高分母 core の正符号 graph の候補と matching 型の係数制約は利用できるが，
局所整数族の存在だけから五 core の接着を結論しない．

優先順位は固定しない．次の局所構造，整数接着の障害，幾何的構成への寄与を見て，
研究としての面白さとインパクトを最大化するように変更する．
